\chapter{2008年福建卷（文）}
\section{单选题}

\begin{problem}
若集合 \(A = \{x \mid x^2 - x < 0\}\)，\(B = \{x \mid 0 < x < 3\}\)，则 \(A \cap B\) 等于（\quad）

\choices
    {\(\{x \mid 0 < x < 1\}\)}
    {\(\{x \mid 0 < x < 3\}\)}
    {\(\{x \mid 1 < x < 3\}\)}
    {\(\varnothing\)}

\begin{answer}
A
\end{answer}

\begin{solution}
由 \(x^2-x=x(x-1)<0\)，得 \(0<x<1\)，所以 \(A=\{x\mid 0<x<1\}\). 又 \(B=\{x\mid 0<x<3\}\)，因此 \(A\cap B=\{x\mid 0<x<1\}\)，故选 A.
\end{solution}
\end{problem}

\begin{problem}
“ \(a = 1\) ”是“直线 \(x + y = 0\) 和直线 \(x - ay = 0\) 互相垂直”的（\quad）

\choices
    {充分而不必要条件}
    {必要而不充分条件}
    {充要条件}
    {既不充分也不必要条件}

\begin{answer}
C
\end{answer}

\begin{solution}
直线 \(x+y=0\) 的一个方向向量为 \((1,-1)\)，直线 \(x-ay=0\) 的一个方向向量为 \((a,1)\). 两直线互相垂直当且仅当 \((1,-1)\cdot(a,1)=a-1=0\)，即 \(a=1\). 因此“\(a=1\)”既是两直线互相垂直的充分条件，也是必要条件，故选 C.
\end{solution}
\end{problem}

\begin{problem}
设 \(\{a_{n}\}\) 是等差数列，若 \(a_{2} = 3\)，\(a_{1} = 13\)，则数列 \(\{a_{n}\}\) 前 8 项的和为（\quad）

\choices
    {128}
    {80}
    {64}
    {56}

\begin{answer}
\(-176\)（按题面数据计算，四个选项均不对应）
\end{answer}

\begin{solution}
设等差数列 \(\{a_n\}\) 的公差为 \(d\). 由 \(a_2=a_1+d\)，得 \(d=a_2-a_1=3-13=-10\). 因此前 8 项的和为 \(S_8=\frac{8}{2}\bigl(2a_1+7d\bigr)=4\bigl(2\times13+7\times(-10)\bigr)=-176\).

所以按原题面计算，答案应为 \(-176\)，但原题给出的四个选项中没有该值.
\end{solution}
\end{problem}

\begin{problem}
函数 \(f(x) = x^3 + \sin x + 1\)（\(x \in \R\)），若 \(f(a) = 2\)，则 \(f(-a)\) 的值为（\quad）

\choices
    {3}
    {0}
    {\(-1\)}
    {\(-2\)}

\begin{answer}
B
\end{answer}

\begin{solution}
由 \(f(a)=2\)，得 \(a^3+\sin a=1\). 又 \(f(-a)=(-a)^3+\sin(-a)+1=-a^3-\sin a+1=0\)，故选 B.
\end{solution}
\end{problem}

\begin{problem}
某一批花生种子，如果每 1 粒发芽的概率为 \(\frac{4}{5}\)，那么播下 3 粒种子恰有 2 粒发芽的概率是（\quad）

\choices
    {\(\frac{12}{125}\)}
    {\(\frac{16}{125}\)}
    {\(\frac{48}{125}\)}
    {\(\frac{96}{125}\)}

\begin{answer}
C
\end{answer}

\begin{solution}
设三粒种子发芽的事件分别为 \(A_1\)，\(A_2\)，\(A_3\). 由独立性，恰有两粒发芽包括三种互不相容的情形，因此
\[
P=P(A_1A_2\overline{A_3})+P(A_1\overline{A_2}A_3)+P(\overline{A_1}A_2A_3)
=\left(\frac45\right)^2\frac15+\frac45\frac15\frac45+\frac15\left(\frac45\right)^2
=3\left(\frac45\right)^2\frac15=\frac{48}{125}
\]
所以选 C.
\end{solution}
\end{problem}

\begin{problem}
如图，在长方体 \(ABCD - A_{1}B_{1}C_{1}D_{1}\) 中，\(AB = BC = 2\)，\(AA_{1} = 1\)，则 \(AC_{1}\) 与平面 \(A_{1}B_{1}C_{1}D_{1}\) 所成角的正弦值为（\quad）

\bitmapfigure[max width=0.42\linewidth]{img/2008/fujian_liberal/q6_fig1.png}

\choices
    {\(\frac{2\sqrt{2}}{3}\)}
    {\(\frac{2}{3}\)}
    {\(\frac{\sqrt{2}}{4}\)}
    {\(\frac{1}{3}\)}

\begin{answer}
D
\end{answer}

\begin{solution}
线段 \(AC_1\) 在平面 \(A_1B_1C_1D_1\) 上的射影为 \(A_1C_1\). 由 \(AB=BC=2\)，得 \(A_1C_1=\sqrt{AB^2+BC^2}=2\sqrt2\). 又 \(AA_1=1\)，所以 \(AC_1=\sqrt{A_1C_1^2+AA_1^2}=\sqrt{8+1}=3\). 设 \(AC_1\) 与该平面所成的角为 \(\theta\)，则 \(\sin\theta=\dfrac{AA_1}{AC_1}=\dfrac13\)，故选 D.
\end{solution}
\end{problem}

\begin{problem}
函数 \(y = \cos x\)（\(x \in \R\)）的图象向左平移 \(\frac{\pi}{2}\) 个单位后，得到函数 \(y = g(x)\) 的图象，则 \(g(x)\) 的解析式为（\quad）

\choices
    {\(-\sin x\)}
    {\(\sin x\)}
    {\(-\cos x\)}
    {\(\cos x\)}

\begin{answer}
A
\end{answer}

\begin{solution}
函数 \(y=\cos x\) 的图象向左平移 \(\frac\pi2\) 个单位，所得函数为 \(g(x)=\cos\left(x+\frac\pi2\right)=-\sin x\)，故选 A.
\end{solution}
\end{problem}

\begin{problem}
在 \(\triangle ABC\) 中，角 \(A\)，\(B\)，\(C\) 的对边分别为 \(a\)，\(b\)，\(c\)，若 \(a^2 + c^2 - b^2 = \sqrt{3}ac\)，则角 \(B\) 的值为（\quad）

\choices
    {\(\frac{\pi}{6}\)}
    {\(\frac{\pi}{3}\)}
    {\(\frac{\pi}{3}\) 或 \(\frac{5\pi}{6}\)}
    {\(\frac{\pi}{3}\) 或 \(\frac{2\pi}{3}\)}

\begin{answer}
A
\end{answer}

\begin{solution}
由余弦定理，\(b^2=a^2+c^2-2ac\cos B\). 结合题设得 \(2ac\cos B=\sqrt3ac\). 因为三角形边长 \(a\)，\(c>0\)，所以 \(\cos B=\frac{\sqrt3}{2}\). 由于 \(0<B<\pi\)，故 \(B=\frac\pi6\)，选 A.
\end{solution}
\end{problem}

\begin{problem}
某班级要从 4 名男生，2 名女生中选派 4 人参加某次社区服务，如果要求至少有 1 名女生，那么不同的选派方案种数为（\quad）

\choices
    {14}
    {24}
    {28}
    {48}

\begin{answer}
A
\end{answer}

\begin{solution}
从 6 人中任取 4 人共有 \(\mathrm C_6^4\) 种方案，其中不含女生的方案只有从 4 名男生中全取，计 \(\mathrm C_4^4\) 种. 因而至少有 1 名女生的方案数为 \(\mathrm C_6^4-\mathrm C_4^4=15-1=14\)，故选 A.
\end{solution}
\end{problem}

\begin{problem}
若实数 \(x\)，\(y\) 满足 \(\begin{cases}
x - y + 1 \leqslant 0,\\
x > 0,\\
y \leqslant 2,
\end{cases}\) 则 \(\frac{y}{x}\) 的取值范围是（\quad）

\choices
    {\((0,2)\)}
    {\((0,2]\)}
    {\((2, +\infty)\)}
    {\([2, +\infty)\)}

\begin{answer}
D
\end{answer}

\begin{solution}
由 \(x-y+1\leqslant0\)，得 \(y\geqslant x+1\). 又 \(x>0\) 且 \(y\leqslant2\)，所以 \(x+1\leqslant2\)，即 \(0<x\leqslant1\). 因此
\[
\frac yx\geqslant\frac{x+1}{x}=1+\frac1x\geqslant2
\]
当 \(r\geqslant2\) 时，取 \(x=\frac1{r-1}\)，\(y=\frac r{r-1}\)，则 \(0<x\leqslant1\)、\(y\leqslant2\)、\(y\geqslant x+1\)，且 \(\frac yx=r\). 所以 \(\frac yx\) 的取值范围为 \([2,+\infty)\)，选 D.
\end{solution}
\end{problem}

\begin{problem}
如果函数 \(y = f(x)\) 的图象如下图，那么导函数 \(y = f'(x)\) 的图象可能是（\quad）

\bitmapfigure[max width=0.88\linewidth]{img/2008/fujian_liberal/q11_fig1.png}

\choices
    {\choicebitmap[width=0.15\paperwidth]{img/2008/fujian_liberal/q11_optA.png}}
    {\choicebitmap[width=0.15\paperwidth]{img/2008/fujian_liberal/q11_optB.png}}
    {\choicebitmap[width=0.15\paperwidth]{img/2008/fujian_liberal/q11_optC.png}}
    {\choicebitmap[width=0.15\paperwidth]{img/2008/fujian_liberal/q11_optD.png}}

\begin{answer}
A
\end{answer}

\begin{solution}
由题图可知，函数 \(f(x)\) 先增后减，再增后减，因此 \(f'(x)\) 的符号应依次为正、负、正、负；在 \(f(x)\) 的三个极值点处，\(f'(x)=0\). 选项 A 的图象具有上述符号变化和零点位置，其他选项至少有一处与 \(f(x)\) 的单调性不符，故选 A.
\end{solution}
\end{problem}

\begin{problem}
双曲线 \(\frac{x^2}{a^2} - \frac{y^2}{b^2} = 1\)（\(a > 0\)，\(b > 0\)）的两个焦点为 \(F_1\)，\(F_2\). 若 \(P\) 为其上一点，且 \(|PF_1| = 2|PF_2|\)，则双曲线离心率的取值范围为（\quad）

\choices
    {\((1,3)\)}
    {\((1,3]\)}
    {\((3, +\infty)\)}
    {\([3, +\infty)\)}

\begin{answer}
B
\end{answer}

\begin{solution}
设双曲线半焦距为 \(c\)，则离心率 \(e=\frac ca\)，且 \(c^2=a^2+b^2\)，所以 \(e>1\). 双曲线的定义给出 \(\left||PF_1|-|PF_2|\right|=2a\). 设较短的焦半径为 \(r\)，则由题设较长的焦半径为 \(2r\)，于是 \(r=2a\)，两焦半径分别为 \(2a\)，\(4a\).

由三角形不等式，焦点间距离 \(2c\) 满足 \(2a\leqslant2c\leqslant6a\)，从而 \(1\leqslant e\leqslant3\). 结合 \(b>0\) 得 \(e>1\)，故 \(1<e\leqslant3\). 反过来，当 \(1<e\leqslant3\) 时，以两焦点为圆心、\(2a\) 和 \(4a\) 为半径的两圆满足交点存在条件 \(2a\leqslant2c\leqslant6a\)，交点即为满足题设距离比的双曲线上一点，因此范围确为 \((1,3]\)，选 B.
\end{solution}
\end{problem}

\section{填空题}

\begin{problem}
\(\left(x + \frac{1}{x}\right)^9\) 展开式中 \(x^2\) 的系数是\fillinblank{}.（用数字作答）

\begin{answer}
\(0\)
\end{answer}

\begin{solution}
按二项式定理，展开式的第 \(k+1\) 项为 \(\mathrm C_9^k x^{9-k}\left(\frac1x\right)^k=\mathrm C_9^k x^{9-2k}\)，其中 \(k=0\)，\(1\)，\(\ldots\)，\(9\). 指数 \(9-2k\) 始终为奇数，不可能等于 2，所以展开式中没有 \(x^2\) 项，其系数为 \(0\).
\end{solution}
\end{problem}

\begin{problem}
若直线 \(3x + 4y + m = 0\) 与圆 \(x^{2} + y^{2} - 2x + 4y + 4 = 0\) 没有公共点，则实数 \(m\) 的取值范围是\fillinblank{}.

\begin{answer}
\((-\infty,0)\cup(10,+\infty)\)
\end{answer}

\begin{solution}
圆的方程可化为 \((x-1)^2+(y+2)^2=1\)，所以圆心为 \((1,-2)\)，半径为 \(1\). 圆心到直线 \(3x+4y+m=0\) 的距离为
\[
\frac{|3\times1+4\times(-2)+m|}{\sqrt{3^2+4^2}}=\frac{|m-5|}{5}
\]
直线与圆没有公共点当且仅当该距离大于半径，即 \(\frac{|m-5|}{5}>1\). 因而 \(|m-5|>5\)，解得 \(m<0\) 或 \(m>10\)，所以取值范围为 \((-\infty,0)\cup(10,+\infty)\).
\end{solution}
\end{problem}

\begin{problem}
若三棱锥的三个侧面两两垂直，且侧棱长均为 \(\sqrt{3}\)，则其外接球的表面积是\fillinblank{}.

\begin{answer}
\(9\pi\)
\end{answer}

\begin{solution}
设三棱锥的顶点为 \(O\)，底面顶点为 \(A\)，\(B\)，\(C\). 由于三个侧面两两垂直，可取 \(OA\)，\(OB\)，\(OC\) 为互相垂直的三条坐标轴，并设 \(A=(\sqrt3,0,0)\)，\(B=(0,\sqrt3,0)\)，\(C=(0,0,\sqrt3)\)，\(O=(0,0,0)\).

外接球球心到 \(O\)，\(A\)，\(B\)，\(C\) 的距离相等. 由对称性球心为 \(Q=\left(\frac{\sqrt3}{2},\frac{\sqrt3}{2},\frac{\sqrt3}{2}\right)\)，其半径满足 \(R=QO=\sqrt{3\left(\frac{\sqrt3}{2}\right)^2}=\frac32\). 因此外接球的表面积为 \(4\pi R^2=4\pi\left(\frac32\right)^2=9\pi\).
\end{solution}
\end{problem}

\begin{problem}
设 \(P\) 是一个数集，且至少含有两个数，若对任意 \(a\)，\(b \in P\)，都有 \(a+b\)，\(a-b\)，\(ab\)，\(\frac{a}{b} \in P\)（除数 \(b \neq 0\)），则称 \(P\) 是一个数域. 例如有理数集 \(\mathbf{Q}\) 是数域. 有下列命题：

\begin{circlelist}
\item 数域必含有 \(0,1\) 两个数；
\item 整数集是数域；
\item 若有理数集 \(\mathbf{Q} \subseteq M\)，则数集 \(M\) 必为数域；
\item 数域必为无限集.
\end{circlelist}

其中正确的命题的序号是\fillinblank{}.（把你认为正确的命题的序号都填上）

\begin{answer}
\(\circled{1}\)，\(\circled{4}\)
\end{answer}

\begin{solution}
判断四个命题如下.

（1）取数域 \(P\) 中两个不同的数 \(a\)，\(b\). 因为 \(a-b\in P\) 且 \(a-b\neq0\)，所以 \(1=\frac{a-b}{a-b}\in P\)，再由封闭性得 \(0=1-1\in P\)，命题正确.

（2）整数集不是数域，因为 \(1\)，\(2\in\Z\)，但 \(\frac12\notin\Z\)，命题错误.

（3）例如取 \(M=\mathbf{Q}\cup\{\sqrt2\}\). 虽然 \(\mathbf{Q}\subseteq M\)，但 \(1+\sqrt2\notin M\)，所以 \(M\) 不一定是数域，命题错误.

（4）数域含有 \(1\)，由加法封闭性含有 \(1\)，\(2\)，\(3\)，\(\ldots\)，这些数两两不同，因此数域必为无限集，命题正确. 所以正确命题的序号为 \(\circled{1}\)，\(\circled{4}\).
\end{solution}
\end{problem}

\section{解答题}

\begin{problem}
已知向量 \(\bs{m} = (\sin A, \cos A)\)，\(\bs{n} = (1, -2)\)，且 \(\bs{m} \cdot \bs{n} = 0\).

\begin{enumerate}
\item 求 \(\tan A\) 的值；
\item 求函数 \(f(x) = \cos 2x + 4\cos A \sin x\)（\(x \in \R\)）的值域.
\end{enumerate}
\end{problem}

\begin{solution}
\begin{enumerate}
\item 由 \(\bs{m} \cdot \bs{n} = 0\)，得
\(\sin A-2\cos A=0\). 若 \(\cos A=0\)，则上式不成立，因此 \(\cos A\neq0\)，从而 \(\tan A=\frac{\sin A}{\cos A}=2\).
\item 由 \(\sin A=2\cos A\) 及 \(\sin ^2A+\cos ^2A=1\)，得
\(5\cos ^2A=1\)，\(\left|\cos A\right|=\frac{1}{\sqrt{5}}\). 令 \(t=\sin x\)，则 \(t\in[-1,1]\)，且 \(f(x)=1-2t^2+4t\cos A=1+2\cos ^2A-2\left(t-\cos A\right)^2\).
因为 \(\left|\cos A\right|\leqslant1\)，所以 \(t=\cos A\) 可以取到，此时 \(f(x)\) 取得最大值
\(1+2\cos ^2A=\frac{7}{5}\).
这是关于 \(t\) 的开口向下的二次函数，所以最小值在 \(t=-1\) 或 \(t=1\) 处取得. 又
\(f\left(\frac{\pi}{2}\right)=-1+4\cos A\)，\(f\left(-\frac{\pi}{2}\right)=-1-4\cos A\)，
故最小值为 \(-1-4\left|\cos A\right|=-1-\frac{4}{\sqrt{5}}\). 因此，函数 \(f(x)\) 的值域为 \(\left[-1-\frac{4}{\sqrt{5}},\frac{7}{5}\right]\).
\end{enumerate}
\end{solution}

\begin{problem}
三人独立破译同一份密码. 已知三人各自破译出密码的概率分别为 \(\frac{1}{5}\)，\(\frac{1}{4}\)，\(\frac{1}{3}\)，且他们是否破译出密码互不影响.

\begin{enumerate}
\item 求恰有二人破译出密码的概率；
\item “密码被破译”与“密码未被破译”的概率哪个大？说明理由.
\end{enumerate}
\end{problem}

\begin{solution}
设第 \(i\) 个人破译出密码为事件 \(A_i\)（\(i=1\)，\(2\)，\(3\)）.

\begin{enumerate}
\item 恰有二人破译出密码包括三种互不相容的情形，因此所求概率为
\[
P=\frac{1}{5}\cdot\frac{1}{4}\cdot\left(1-\frac{1}{3}\right)+\frac{1}{5}\cdot\left(1-\frac{1}{4}\right)\cdot\frac{1}{3}+\left(1-\frac{1}{5}\right)\cdot\frac{1}{4}\cdot\frac{1}{3}=\frac{1}{30}+\frac{1}{20}+\frac{1}{15}=\frac{3}{20}
\]
\item 三人都未破译出密码的概率为 \(P(\text{密码未被破译})=\left(1-\frac{1}{5}\right)\left(1-\frac{1}{4}\right)\left(1-\frac{1}{3}\right)=\frac{4}{5}\cdot\frac{3}{4}\cdot\frac{2}{3}=\frac{2}{5}\). 所以 \(P(\text{密码被破译})=1-\frac{2}{5}=\frac{3}{5}>\frac{2}{5}\)，
故“密码被破译”的概率较大.
\end{enumerate}
\end{solution}

\begin{problem}
如图，在四棱锥 \(P - ABCD\) 中，侧面 \(PAD \perp\) 底面 \(ABCD\)，侧棱 \(PA = PD = \sqrt{2}\)，底面 \(ABCD\) 为直角梯形，其中 \(BC \parallel AD\)，\(AB \perp AD\)，\(AD = 2AB = 2BC = 2\)，\(O\) 为 \(AD\) 中点.

\begin{enumerate}
\item 求证：\(PO \perp\) 平面 \(ABCD\)；
\item 求异面直线 \(PB\) 与 \(CD\) 所成角的余弦值；
\item 求点 \(A\) 到平面 \(PCD\) 的距离.
\end{enumerate}

\bitmapfigure[max width=0.42\linewidth]{img/2008/fujian_liberal/q19_fig1.png}
\end{problem}

\begin{solution}
\begin{enumerate}
\item 在等腰三角形 \(PAD\) 中，\(PA=PD\)，且 \(O\) 是 \(AD\) 的中点，所以由 \(\triangle PAO\cong\triangle PDO\) 得 \(PO\perp AD\). 又平面 \(PAD\perp\) 平面 \(ABCD\)，且两平面的交线为 \(AD\)，因此 \(PO\perp\) 平面 \(ABCD\).
\item 建立空间直角坐标系，使平面 \(ABCD\) 为 \(z=0\)，取
\(A=(0,0,0)\)，\(D=(2,0,0)\)，\(B=(0,1,0)\)，\(C=(1,1,0)\).
由 \(O=(1,0,0)\)，\(AO=1\)，\(PA=\sqrt{2}\)，以及 \(PO\perp\) 平面 \(ABCD\)，得 \(PO=1\)，可取 \(P=(1,0,1)\).

于是 \(\overrightarrow{PB}=(-1,1,-1)\)，\(\overrightarrow{CD}=(1,-1,0)\).
异面直线所成的角取锐角，所以其余弦值为
\[
\frac{\left|\overrightarrow{PB}\cdot\overrightarrow{CD}\right|}{\left|\overrightarrow{PB}\right|\left|\overrightarrow{CD}\right|}
=\frac{2}{\sqrt{3}\sqrt{2}}=\frac{\sqrt{6}}{3}
\]
\item 由上面的坐标，取
\(\overrightarrow{PC}=(0,1,-1)\)，\(\overrightarrow{PD}=(1,0,-1)\).
向量 \((1,1,1)\) 同时垂直于这两个向量，因此平面 \(PCD\) 的方程为
\((x-1)+y+(z-1)=0\)，即 \(x+y+z-2=0\).
所以点 \(A\) 到平面 \(PCD\) 的距离为
\[
\frac{\left|0+0+0-2\right|}{\sqrt{1^2+1^2+1^2}}=\frac{2}{\sqrt{3}}=\frac{2\sqrt{3}}{3}
\]
\end{enumerate}
\end{solution}

\begin{problem}
已知 \(\{a_{n}\}\) 是正数组成的数列，\(a_{1} = 1\)，且点 \((\sqrt{a_{n}}, a_{n+1})\)（\(n \in \mathbf{N}^{*}\)）在函数 \(y = x^{2} + 1\) 的图象上.

\begin{enumerate}
\item 求数列 \(\{a_{n}\}\) 的通项公式；
\item 若数列 \(\{b_{n}\}\) 满足 \(b_{1}=1\)，\(b_{n+1}=b_{n}+2^{a_{n}}\)，求证：\(b_{n}\cdot b_{n+2}<b_{n+1}^{2}\).
\end{enumerate}
\end{problem}

\begin{solution}
\begin{enumerate}
\item 因为点 \(\left(\sqrt{a_n},a_{n+1}\right)\) 在函数 \(y=x^2+1\) 的图象上，所以
\(a_{n+1}=\left(\sqrt{a_n}\right)^2+1=a_n+1\). 又 \(a_1=1\)，由此递推可得 \(a_n=1+(n-1)=n\). 故数列 \(\{a_n\}\) 的通项公式为 \(a_n=n\).
\item 由 \(a_n=n\)，递推关系化为
\(b_{n+1}=b_n+2^n\). 由 \(b_1=1=2^1-1\)，并且当 \(b_n=2^n-1\) 时 \(b_{n+1}=2^n-1+2^n=2^{n+1}-1\)，所以由数学归纳法得 \(b_n=2^n-1\).

于是
\[
b_{n+1}^2-b_nb_{n+2}=\left(2^{n+1}-1\right)^2-\left(2^n-1\right)\left(2^{n+2}-1\right)=2^n>0
\]
因此 \(b_n\cdot b_{n+2}<b_{n+1}^2\).
\end{enumerate}
\end{solution}

\begin{problem}
已知函数 \(f(x) = x^3 + mx^2 + nx - 2\) 的图象过点 \((-1, -6)\)，且函数 \(g(x) = f'(x) + 6x\) 的图象关于 \(y\) 轴对称.

\begin{enumerate}
\item 求 \(m\)，\(n\) 的值及函数 \(y = f(x)\) 的单调区间；
\item 若 \(a > 0\)，求函数 \(y = f(x)\) 在区间 \((a - 1, a + 1)\) 内的极值.
\end{enumerate}
\end{problem}

\begin{solution}
\begin{enumerate}
\item 由函数图象过点 \((-1,-6)\)，得 \(-1+m-n-2=-6\)，即 \(m-n=-3\).
又 \(g(x)=f'(x)+6x=3x^2+(2m+6)x+n\).
由于 \(g(x)\) 的图象关于 \(y\) 轴对称，所以 \(g(x)\) 是偶函数，故 \(2m+6=0\)，从而
\(m=-3\)，\(n=0\).
因此
\(f(x)=x^3-3x^2-2\)，\(f'(x)=3x^2-6x=3x(x-2)\).
当 \(x<0\) 或 \(x>2\) 时，\(f'(x)>0\)；当 \(0<x<2\) 时，\(f'(x)<0\). 所以函数 \(f(x)\) 在 \((-\infty,0)\) 和 \((2,+\infty)\) 上单调递增，在 \((0,2)\) 上单调递减.
\item 由上面的单调性，\(x=0\) 是极大值点，且 \(f(0)=-2\)；\(x=2\) 是极小值点，且 \(f(2)=-6\).

又因为 \(a>0\)，所以 \(0\in(a-1,a+1)\) 当且仅当 \(0<a<1\)，\(2\in(a-1,a+1)\) 当且仅当 \(1<a<3\).
因此，当 \(0<a<1\) 时，函数在该区间内取得极大值 \(-2\)；当 \(1<a<3\) 时，函数在该区间内取得极小值 \(-6\)；当 \(a=1\) 或 \(a\geqslant3\) 时，区间内没有极值.
\end{enumerate}
\end{solution}

\begin{problem}
如图，椭圆 \(C: \frac{x^2}{a^2} + \frac{y^2}{b^2} = 1\)（\(a > b > 0\)）的一个焦点为 \(F(1,0)\)，且过点 \((2,0)\).

\begin{enumerate}
\item 求椭圆 \(C\) 的方程；
\item 若 \(AB\) 为垂直于 \(x\) 轴的动弦，直线 \(l: x = 4\) 与 \(x\) 轴交于点 \(N\)，直线 \(AF\) 与 \(BN\) 交于点 \(M\).

\begin{enumerate}
\item 求证：点 \(M\) 恒在椭圆 \(C\) 上；
\item 求 \(\triangle AMN\) 面积的最大值.
\end{enumerate}

\bitmapfigure[max width=0.42\linewidth]{img/2008/fujian_liberal/q22_fig1.png}
\end{enumerate}
\end{problem}

\begin{solution}
\begin{enumerate}
\item 设椭圆的半焦距为 \(c\)，则 \(c=1\)，且 \(a^2-b^2=c^2=1\). 因为点 \((2,0)\) 在椭圆上，所以 \(\frac{2^2}{a^2}=1\). 由 \(a>0\)，得 \(a=2\)，从而 \(b^2=a^2-1=3\). 所以椭圆 \(C\) 的方程为
\[
\frac{x^2}{4}+\frac{y^2}{3}=1
\]
\item
\begin{enumerate}
\item 由 \(AB\) 垂直于 \(x\) 轴，设 \(A=(t,s)\)，\(B=(t,-s)\)，\(-2<t<2\)，\(s>0\). 因为 \(A\) 在椭圆上，所以 \(\frac{t^2}{4}+\frac{s^2}{3}=1\)，从而 \(s^2=\frac{3}{4}\left(4-t^2\right)\).

又 \(F=(1,0)\)，\(N=(4,0)\). 设 \(M=AF\cap BN\)，并令 \(M=A+\lambda(F-A)=B+\mu(N-B)\). 比较两式的纵坐标，得 \(s(1-\lambda)=s(\mu-1)\)，所以 \(\mu=2-\lambda\). 再比较横坐标，得 \(t+\lambda(1-t)=t+(2-\lambda)(4-t)\)，即 \(\lambda(5-2t)=2(4-t)\)，所以 \(\lambda=\frac{2(4-t)}{5-2t}\). 因此 \(M=\left(\frac{8-5t}{5-2t},-\frac{3s}{5-2t}\right)\).

代入椭圆方程，得
\[
\frac{x_M^2}{4}+\frac{y_M^2}{3}=\frac{(8-5t)^2+9(4-t^2)}{4(5-2t)^2}=\frac{100-80t+16t^2}{4(5-2t)^2}=1
\]
所以点 \(M\) 恒在椭圆 \(C\) 上.

\item 设 \(S=S_{\triangle AMN}\)，则 \(M-A=\left(\frac{2(t-1)(t-4)}{5-2t},-\frac{2s(4-t)}{5-2t}\right)\)，\(N-A=(4-t,-s)\). 因此
\[
S=\frac{1}{2}\left|\det(M-A,N-A)\right|=\frac{3s(4-t)}{5-2t}=\frac{3\sqrt{3}}{2}\frac{(4-t)\sqrt{4-t^2}}{5-2t}
\]
在 \(-2<t<2\) 内，最大化 \(S\) 等价于最大化 \(H(t)=\frac{(4-t)^2(4-t^2)}{(5-2t)^2}\). 对数求导，得
\[
\frac{H'(t)}{H(t)}=-\frac{2}{4-t}-\frac{2t}{4-t^2}+\frac{4}{5-2t}
\]
令上式等于零并通分化简，得
\[
0=\frac{-4(t-1)(t^2-4t+6)}{(4-t)(4-t^2)(5-2t)}
\]
因为 \(-2<t<2\)，分母为正；又 \(t^2-4t+6=(t-2)^2+2>0\)，所以唯一的临界点为 \(t=1\). 且当 \(t\) 趋近于 \(-2\) 或 \(2\) 时，\(S\) 趋近于 \(0\)，故最大值在 \(t=1\) 时取得. 此时 \(s=\frac{3}{2}\)，所以 \(S_{\max}=\frac{3\cdot\frac{3}{2}\cdot3}{3}=\frac{9}{2}\). 故 \(\triangle AMN\) 面积的最大值为 \(\frac{9}{2}\).
\end{enumerate}
\end{enumerate}
\end{solution}
